%!TEX root = ./../../springer-mono-format/book.tex
%!TEX root = ./../../thesis-trier.tex

\index{OI-automaton|see{automaton with one-point singular idempotents}} \IN{automaton with one-point singular idempotents}
\index{SI-automaton|see{automaton with simple idempotents}} \IN{automaton with simple idempotents}
Let $\semiautomat{A}{Q}{\Sigma}{\delta}$.
We call $\mathcal A$ an \emph{automaton with one-point singular idempotents (OI-automaton)}
if every letter is either permutational or induces a one-point singular idempotent.
We call $\mathcal A$ an \emph{automaton with simple idempotents (SI-automaton)}
if every letter is either permutational or a simple idempotent.
Clearly, every SI-automaton is an OI-automaton.

For example, the \Cerny automata from
Section~\ref{intro:cerny_rt} (Figure~\ref{intro:fig:cerny})
are SI-automata.
Rystsov~\cite{Rystsov2000} introduced SI-automata and proved that
every synchronizing SI-automaton with $n$ states admits a synchronizing
word of length at most $2(n-1)^2$.
In \cref{sync_max:graphs:sec:extend-bound}, we extend this bound to OI-automata.

For an OI-automaton  $\semiautomat{A}{Q}{\Sigma}{\delta}$, we set \label{def:subalphabets_SI}
\begin{align*} 
	\Sigma_I & = \{ a \in \Sigma \mid a \text{ induces a one-point singular idempotent } \}, \\
	\Sigma_G & = \{ a \in \Sigma \mid a \text{ is a permutational letter } \}.
\end{align*} 
Then $\Sigma = \Sigma_I \cup \Sigma_G$. 

First, in
\cref{prop:structure_simp_idem_binary},
we classify synchronizing binary SI-automata, a result originally
due to Martyugin~\cite[Prop.~6.2]{Mar2009b}.
\Cref{prelim:prop:binary-SI-sync-words} gives synchronizing words for
binary SI-automata.
\Cref{lem:each_sync_word_contains_a} establishes a lower bound on the rank
of a reachable subset, and
\cref{cor:each_sync_word_contains_a} establishes a lower bound on the
length of a word of a given rank.
\Cref{prelim:prop:SI-without-perm:reach-bound} determines the length of shortest
words reaching a given subset in SI-automata.

\begin{proposition}[{Martyugin~\cite[Prop.~6.2]{Mar2009b}}] \label{prop:structure_simp_idem_binary}
	Let $\semiautomat{A}{Q}{\{a,b\}}{\delta}$ be an $n$-state binary synchronizing SI-automaton.
	Then we have $n \le 2$ or one of the following holds:
	\begin{enumerate}
		\item $n = 3$ and $\mathcal A$ is, up to renaming of the letters, one of the automata in \cref{prelim:fig:three-state-sync-binary-SI-a,prelim:fig:three-state-sync-binary-SI-b}, or
		\item one letter, say $a$, is a simple idempotent with $\delta(s,a)=t$ for distinct
		$s,t\in Q$, the letter $b$ permutes the states with
		$
		Q=\{\delta(s,b^i)\mid i\ge0\}\cup\{t\},
		$
		and either
		\begin{enumerate}
			\item $t$ is a trap state (see \cref{prelim:fig:binarySI-sync-trap}), or
			\item $\delta(s,b^d)=t$ for some $0<d<n$ with $\gcd(d,n)=1$
			(see \cref{prelim:fig:binarySI-sync-SCC}).\footnote{In~\cite[Prop.~6.2]{Mar2009b} a 
				different naming of the states is used,
				with a parameter $2\le p\le n$ such that $p=n-d+1$. There appears to be a minor
				mistake in~\cite[Prop.~6.2]{Mar2009b}. It is claimed that the 
				automaton is not synchronizing whenever
				$\gcd(n,p)>1$. However, if $p=n$ (which results in the \Cerny automaton),
				then $\gcd(n,p)=n$. The correct condition is
				$\gcd(n,n-p+1)>1$.}
		\end{enumerate}
	\end{enumerate}
\end{proposition}
\begin{proof}
	Clearly, one letter, say~$a$, must be a simple idempotent.
		
	First, suppose~$a$ and~$b$ are both simple idempotents.
	If $n \ge 4$, then, since each letter moves only one state,
	there exist two distinct states $p,q \in Q$
	that are fixed by both letters.
	Consequently, every word fixes them, and so
	no word synchronizes the automaton.
	Thus, we have $n \le 3$.
	It is clear that for $n \le 2$, $\mathcal A$
	is always synchronizing.
	For $n = 3$, 
	a simple case analysis yields that~$\mathcal A$
	is synchronizing if and only if it equals,
	up to renaming of the letters, one 
	of the automata from \cref{prelim:fig:three-state-sync-binary-SI-a,prelim:fig:three-state-sync-binary-SI-b}.
	
	Next, let~$b$ be a permutational letter. Let $s,t \in Q$ be distinct with $\delta(s, a) = t$.
	Set $S = \{\delta(s, b^i) \mid i \ge 0\}$
	and $T = \{\delta(t, b^i) \mid i \ge 0\}$.
	As $\delta(S \cup T, x) \subseteq S \cup T$
	and $\delta(Q \setminus (S \cup T), x) \subseteq Q \setminus (S \cup T)$ for $x \in \{a,b\}$,
	we have $Q = S\cup T$.	
	If $s \notin T$, then $\delta(T, x) = T$
	for all $x \in \Sigma$, which gives
	$T = \{t\}$ and that~$t$ is a trap state.
	Otherwise, $s \in T$ implies $S = T$. Then
	$\delta(s, b^d) = t$ for some 
	$0 < d < n$. Set $r = \gcd(d, n)$
	and, for $p, q \in Q$, write $p \sim q$
	if there exists $i \ge 0$ such that $\delta(p, b^{ir}) = q$.
	For all $x \in \{a,b\}$ and $p,q \in Q$ we have
	\[
	p \sim q \Leftrightarrow \delta(p, x) \sim \delta(q, x).
	\]
	So $|\delta(\{p,q\}, u)| = 2$ for all $p, q \in Q$ with $p \not\sim q$.
	Hence, as $\mathcal A$ is synchronizing, all states must be equivalent
	w.r.t.~$\sim$, which implies $r = 1$.	
%	
%		Next, suppose~$b$ permutes the states in a single cycle.
%		Then we can assume $Q = [n]$ and $b$ adds one modulo $n$
%		to each state. For all states $p,q \in Q$
%		we have that 
%		\[p \equiv q \pmod{\gcd(d,n)} \Leftrightarrow 
%		\applytrans{\delta}{p}{x} \equiv \applytrans{\delta}{q}{x} \pmod{\gcd(d,n)}
%		\]
%		for each $x \in \Sigma$.
%		So, it is not possible to map two states in different residue classes
%		to the same state and hence, for $\gcd(d,n) > 1$,
%		the automaton is not synchronizable. 
\end{proof}

%!TEX root = ./springer-mono-format/book.tex
%!TEX root = ./thesis-trier.tex


\begin{figure}[htb]
	\centering
	
	%--------------------------------------------------------------- row 1
	\begin{subfigure}[t]{.48\textwidth}
		\centering
		\scalebox{0.8}{%
			\begin{tikzpicture}[
				baseline=(current bounding box.north),
				node distance=2.4em
				]
				\node[state]              (1) {};
				\node[state, right=of 1]  (2) {};
				\node[state, right=of 2]  (3) {};
				
				\path
				(1) edge node[above] {$a$} (2)
				(2) edge node[above] {$b$} (3);
			\end{tikzpicture}%
		}
		\caption{Three states, with a trap state.}
		\label{prelim:fig:three-state-sync-binary-SI-a}
	\end{subfigure}\hfill
	\begin{subfigure}[t]{.48\textwidth}
		\centering
		\scalebox{0.8}{%
			\begin{tikzpicture}[
				baseline=(current bounding box.north),
				node distance=2.4em
				]
				\node[state]              (1) {};
				\node[state, left=of 1]   (2) {};
				\node[state, right=of 1]  (3) {};
				
				\path
				(2) edge node[above] {$a$} (1)
				(3) edge node[above] {$b$} (1);
			\end{tikzpicture}%
		}
		\caption{Three states, both letters merging.}
		\label{prelim:fig:three-state-sync-binary-SI-b}
	\end{subfigure}
	
	\medskip
	
	%--------------------------------------------------------------- row 2
	\begin{subfigure}[t]{.48\textwidth}
		\centering
		\scalebox{0.8}{%
			\begin{tikzpicture}[
				%baseline=(current bounding box.north),
				node distance=4.5em
				]
				\def\R{7em}
				\useasboundingbox (-7.8em,-8.0em) rectangle (11.6em,8.0em);
				
				\foreach \name/\angle/\text in
				{A1/0/s,A2/45/,A3/90/,A4/135/,A5/180/,A6/225/,A8/315/}
				{
					\node[state] (\name) at (\angle:\R) {$\text$};
				}
				
				\node[state,right of=A1] (t) {$t$};
				
				
				
				\path
				(A6) edge node[auto] {$b$} (A5)
				(A5) edge node[auto] {$b$} (A4)
				(A4) edge node[auto] {$b$} (A3)
				(A3) edge node[auto] {$b$} (A2)
				(A2) edge node[auto] {$b$} (A1)
				(A1) edge node[auto] {$b$} (A8)
				(A1) edge node[above] {$a$} (t);
				%(t)  edge[loop above] node {$a,b$} (t);
				
				\node at (270:\R) {\large$\cdots$};
			\end{tikzpicture}%
		}
		\caption{Binary synchronizing SI-automaton with a trap state.}
		\label{prelim:fig:binarySI-sync-trap}
	\end{subfigure}\hfill
	\begin{subfigure}[t]{.48\textwidth}
		\centering
		\scalebox{0.8}{%
			\begin{tikzpicture}[
				%baseline=(current bounding box.north),
				]
				\useasboundingbox (-7.8em,-8.0em) rectangle (11.6em,8.0em);
				\def\R{7em}
				
				\foreach \name/\angle/\text in
				{A1/0/,A2/45/,A3/90/,A4/135/,A5/180/t,A6/225/,A7/270/,A8/315/s}
				{
					\node[state] (\name) at (\angle:\R) {$\text$};
				}			
				
				\path
				(A1) edge node[auto] {$b$} (A8)
				(A7) edge node[auto] {$b$} (A6)
				(A6) edge node[auto] {$b$} (A5)
				(A5) edge node[auto] {$b$} (A4)
				(A3) edge node[auto] {$b$} (A2)
				(A2) edge node[auto] {$b$} (A1);
				
                \draw (A8) edge [auto] node {$a$} (A5);
				
				\node[rotate=22.5] at (292.5:\R) {\large$\cdots$};
				\node[rotate=202.5] at (112.5:\R) {\large$\cdots$};
			\end{tikzpicture}%
		}
		\caption{Strongly connected binary synchronizing SI-automaton where
			$\applytrans{\delta}{s}{b^d}=t$.}
		\label{prelim:fig:binarySI-sync-SCC}
	\end{subfigure}
	
	\caption[Synchronizing binary SI-automata]{Synchronizing binary SI-automata (self-loops omitted)
		as listed in \cref{prop:structure_simp_idem_binary}.
		Note that, in \cref{prelim:fig:binarySI-sync-SCC}, for $d=1$ we obtain the \Cerny family.}
	\label{prelim:fig:three-state-sync-binary-SI}
	\label{fig:theorem}
\end{figure}

Next, we give synchronizing words for each of the two types of automata
from \cref{prop:structure_simp_idem_binary}.

\begin{proposition}	Let $n \ge 2$.
	\label{prelim:prop:binary-SI-sync-words}
	\begin{enumerate}
		\item \label{prelim:prop:binary-SI-sync-words:1} The word $(ab)^{n-2}a$ synchronizes the $n$-state SI-automaton 
		from \cref{prelim:fig:binarySI-sync-trap}.
		
		\item \label{prelim:prop:binary-SI-sync-words:2} The word $(ab^{n-d})^{n-2}a$ synchronizes the $n$-state SI-automaton 
		from \cref{prelim:fig:binarySI-sync-SCC}.
	\end{enumerate}
\end{proposition} 
\begin{proof} We use the notation from \cref{prop:structure_simp_idem_binary}.
	
\ref{prelim:prop:binary-SI-sync-words:1}: 
For $i \in \{1,\ldots,n-1\}$ we have
\[
t = \delta(s, b^i (ab)^{(n - 1) - i}a) = \delta(s, b^i (ab)^{n-2}a).
\]
Since the state set is $\{ s, \delta(s, b), \ldots, \delta(s, b^{n-2})\} \cup \{t\}$ 
and~$t$ is a trap state, we find $\delta(Q, (ab)^{n-2}a) = \{t\}$.

\ref{prelim:prop:binary-SI-sync-words:2}: We can assume $n \ge 3$.
For $i \in \{0,1,\ldots,n-1\}$ set
$Q_i = \{\, s, \delta(s,b^{d}), \ldots, \delta(s,b^{id}) \,\}$.
Since $d$ and~$n$ are coprime, the smallest $x > 0$ with $xd \equiv 0 \pmod{n}$
is $x = n$. Hence the states $\delta(s,b^{jd})$ for $j \in \{0,1,\ldots,n-1\}$
are pairwise distinct, so $|Q_i| = i+1$ and $Q_{n-1} = Q$.

We show by induction on~$i$ that $Q_i \subseteq \delta^{-1}(s, (ab^{n-d})^{i})$.
This is easy to see for $i = 1$. Next, suppose $i > 1$.
By induction hypothesis, we have $Q_{i-1} \subseteq \delta^{-1}(s, (ab^{n-d})^{i-1})$.
For $j \in \{0,1,\ldots,i-1\}$, since $d$ and $n$ are coprime, we have $(j+1)d \not\equiv 0 \pmod{n}$
and thus $\delta(s, b^{(j+1)d}) \ne s$.
This gives
$
\delta(s, b^{(j+1)d} a b^{n-d}) = \delta(s, b^{jd}).
$
With $\delta(s, ab^{n-d}) = s$ 
we find $\delta(Q_i, ab^{n-d}) \subseteq Q_{i-1}$.
Thus $Q_i \subseteq \delta^{-1}(s, (ab^{n-d})^{i})$. This concludes our induction argument.

Now, set $q = \delta(s, b^{(n-1)d})$. We have $\delta(q, (ab^{n-d})^{n-2}) = \delta(s, b^d)$
and so
\[
\delta(Q, (ab^{n-d})^{n-2}) = \delta(Q_{n-2} \cup \{q\}, (ab^{n-d})^{n-2})
= \{s, \delta(s, b^d) \}.
\]
Thus, $\delta(Q, (ab^{n-d})^{n-2}a) = \{ \delta(s, b^d)\}$.
%
%
%
%First, we show that for $i \in \{0,1,\ldots,n-2\}$ we have
%\[
%\delta(Q, (ab^{n-d})^i) = \{ s, \delta(s, b^d), \ldots, \delta(s, b^{(n - (i+1))d}) \}
%\]
%Since $n$ and $d$ are coprime, the smallest $x > 0$ such that $xd \equiv 0 \pmod{n}$
%is $x = n$
%and so $Q = \{ s, \delta(s, b^d), \ldots, \delta(s, b^{(n - 1)d}) \}$.
%This gives the claim for $i = 0$.
%Next, suppose for $i \in \{0,1,\ldots,n-3\}$ we have
%\[
%\delta(Q, (ab^{n-d})^i) = \{ s, \delta(s, b^d), \ldots, \delta(s, b^{(n - (i+1))d}) \}
%\]
%Then
%$
%\delta(Q, (ab^{n-d})^ia) = \{ \delta(s, b^d), \ldots, \delta(s, b^{(n - (i+1))d}) \}
%$
%and so 
%\[ 
%\delta(Q, (ab^{n-d})^{i+1}) = \{ s, \delta(s, b^d), \ldots, \delta(s, b^{(n - (i + 2))d}) \}.
%\]
%By induction, our first claim holds. Then 
%we find $\delta(Q, (ab^{n-d})^{n - 2})a) = \delta(\{ s, \delta(s, b^d)\}, a) = \{\delta(s, b^d)\}$
%and so the word $(ab^{n-d})^{n - 2})a$ synchronizes the automaton.
%
%Note that $\delta(s, b^{j(n - d)}) = \delta(s, b^{(n - j)d})$ for $j \in \{0,1,\ldots,n-1\}$
%and so for $i \in \{0,1,\ldots,n-1\}$ we get
%\[
%\{ s, \delta(s, b^d), \ldots, \delta(s, b^{(n - (i+1))d}) \}
%= Q \setminus \{\delta(s, b^{i(n-d)}), \ldots, \delta(s, b^{n-d}) \}. 
%\]
%We also have
%$
%\delta^{-1}(s, (ab^{n-d})^{i}) = \{ s, \delta(s, b^d), \ldots, \delta(s, b^{id}) \}
%$
%for $i \in \{0, 1,\ldots,n-1\}$.
%
%For an alternative proof, consider the automaton from \cref{prelim:fig:binarySI-sync-SCC} 
%and define an automaton~$\mathcal B$ with the same state 
%over two letters $a$ and $c$ where $c$ induces the same transformation as $b^d$ on the states.
%As $d$ and $n$ are coprime,
%the resulting automaton is isomorphic to the \v{C}ern\'y automaton. So, from the method
%of proof of~\cref{intro:cerny_rt}
%we get the synchronizing word $(ac^{(n-1)})^{n-2}a$ for $\mathcal B$,
%which yields the synchronizing word $(ab^{d(n-1)})^{n-2}a$ for the automaton from \cref{prelim:fig:binarySI-sync-SCC}.
%As $dn - d \equiv n - d \pmod{n}$
%the words $b^{d(n-1)}$ and $b^{n - d}$
%induce the same permutation on the states. Hence,
%the word $(ab^{n - d})^{n-2}a$ synchronizes this automaton as well.	
\end{proof}

Further synchronizing words can be constructed. For example,
for the automata from \cref{prelim:fig:binarySI-sync-trap},
all words in
	\[
\Sigma^*a^+(b((ba^*)^{n-1})^*)a^+)^{n-2}\Sigma^*
\]
synchronize it, or $abbabbabbabba$ synchronizes this automaton for $n = 6$.
For the automata from \cref{prelim:fig:binarySI-sync-SCC} 
all words in 
$
\Sigma^*a^+((ba^*)^{n-p}(ba^*)^n a^+)^{n-2}\Sigma^*
$
synchronize it.

\begin{lemma} \label{lem:each_sync_word_contains_a}
	Let $\semiautomat{A}{Q}{\Sigma}{\delta}$ be an SI-automaton.
	Then for $S \subseteq Q$ and $u \in \Sigma^*$
	we have $|S| - \sum_{a \in \Sigma_I} |u|_a \le |\delta(S, u)|$.
\end{lemma}
\begin{proof}
	Clearly, the claim holds for $u = \varepsilon$.
	Next, let $u = wx$ with $x \in \Sigma$ and $w \in \Sigma^*$.
	Set $T = \delta(S, w)$. Inductively, we have
	\begin{equation} \label{prelim:SI:reach-lower-bound}
		|S| - \sum_{a \in \Sigma_I} |w|_a \le |T|.		
	\end{equation}
	For $x \in \Sigma_I$,
	since $x$ maps only two states to one state,
	we have $|T| - 1 \le |\delta(T,x)|$.
	Thus, with $\delta(S, u) = \delta(T,x)$ and \cref{prelim:SI:reach-lower-bound}, we find
	\[
	|S| - \sum_{a \in \Sigma_I} |u|_a
	= |S| - \left( \sum_{a \in \Sigma_I} |w|_a \right) - 1
	\le |T| - 1
	\le |\delta(S, u)|.
	\]
	For $x \in \Sigma_G$,
	we have $|\delta(T, x)| = |T|$
	and $\sum_{a \in \Sigma_I} |u|_a = \sum_{a \in \Sigma_I} |w|_a$
	and substituting these values into \cref{prelim:SI:reach-lower-bound}
	we get the claim.
\end{proof}

\begin{corollary} \label{cor:each_sync_word_contains_a}
	In an $n$-state SI-automaton  $\semiautomat{A}{Q}{\Sigma}{\delta}$ for a
	word $u$ of rank~$m$
	we have $n - m \le |u|$. 
\end{corollary}
\begin{proof}
	By \cref{lem:each_sync_word_contains_a} we have $n - \sum_{a \in \Sigma_I} |u|_a \le m$.
	This gives $n - m \le \sum_{a \in \Sigma_I} |u|_a \le |u|$.	
\end{proof}

Lastly, we note a bound on shortest words reaching a given subset of states
in SI-automata without permutational letters.

\begin{proposition} \label{prelim:prop:SI-without-perm:reach-bound}
	Let $\semiautomat{A}{Q}{\Sigma}{\delta}$ be an SI-automaton without permutational letters.
	Then for a shortest word $u \in \Sigma^*$ such that $\delta(Q, u) = S$
	we have $|u| = |Q| - |S|$.
\end{proposition}
\begin{proof}	
	We can assume $S \ne Q$.
	Let $u_1, \ldots, u_m \in \Sigma$ and $p_1, \ldots, p_m \in Q$
	be such that $u = u_1 \cdots u_m$ and $\excl{u_i} = \{p_i\}$
	for $i \in \{1,\ldots,m\}$.
	Next, by induction, we show that for $i \in \{1,\ldots,m\}$ we have
	\begin{enumerate}	
		\item $\delta(Q, u_1 \cdots u_i) = Q \setminus \{ p_1, \ldots, p_i \}$ and
		
		\item the states $p_1, \ldots, p_i$ are pairwise distinct.
	\end{enumerate}	
	The claim is clear for $i = 1$. Now let $i > 1$.
	By the minimal choice of~$u$,
	we have $p_i \in \delta(Q, u_1 \cdots u_{i-1})$
	since otherwise the letter $u_i$ fixes every state in $\delta(Q, u_1 \cdots u_{i-1})$
	and thus can be omitted.
	Since, inductively, $\delta(Q, u_1 \cdots u_{i-1}) = Q \setminus \{p_1, \ldots, p_{i-1}\}$,
	we find $\delta(Q, u_1 \cdots u_i) = Q \setminus \{ p_1, \ldots, p_i \}$
	and $p_i \notin \{p_1, \ldots, p_{i-1}\}$.
	Since $p_1, \ldots, p_{i-1}$ are pairwise distinct by induction hypothesis,
	the states $p_1, \ldots, p_i$ are pairwise distinct.
	
	Finally, for $i = m$ we find $|S| = |\delta(Q, u)| = |Q| - m$.
	Thus $m = |Q| - |S|$.	
\end{proof}